\chapter{Description of Quadrotor UAV System}
\label{ch:sys_desc}

\section{Physics of Quadrotor UAVs}
A quadrotor is a rotary-wing Unmanned Aerial Vehicle (UAV) that achieves lift, thrust, and attitude control through the aerodynamic forces generated by four independent, fixed-pitch propellers mounted in a symmetric cross configuration \cite{ref_39}. Unlike traditional helicopters, which rely on complex mechanical swashplates to dynamically alter the blade pitch, a quadrotor regulates its entire flight envelope solely by modulating the angular velocities of its four electric motors \cite{ref_41}.

The mechanical symmetry of the quadrotor allows its operation to be governed by relatively straightforward physical principles. The four rotors are arranged in two counter-rotating pairs, as depicted schematically in Figure \ref{fig:quadcopter_nwu_frame}:
\begin{itemize}
    \item \textbf{Rotors 1 and 3:} Spin in a counter-clockwise (CCW) direction, generating a downward aerodynamic thrust and an upward reactive torque acting in the clockwise (CW) direction on the frame.
    \item \textbf{Rotors 2 and 4:} Spin in a clockwise (CW) direction, generating a downward aerodynamic thrust and an upward reactive torque acting in the counter-clockwise (CCW) direction on the frame.
\end{itemize}

\begin{figure}[h!]
\centering
\fbox{
  \begin{minipage}{0.8\textwidth}
    \centering
    \vspace{1.5cm}
    \textbf{}
    \vspace{1.5cm}
  \end{minipage}
}
\caption{Schematic of the quadrotor frame configuration with North-West-Up (NWU) coordinate convention, denoting the positive directions of roll ($\phi$), pitch ($\theta$), and yaw ($\psi$) angles.}
\label{fig:quadcopter_nwu_frame}
\end{figure}

The cancellation of net aerodynamic torque is achieved during balanced hover when all four rotors rotate at the exact same angular velocity. In this state, the sum of CW torques matches the sum of CCW torques, keeping the heading (yaw) of the quadrotor constant. Operational maneuvers are achieved through differential thrust schemes across specific rotor pairs \cite{ref_42, ref_49}:
\begin{enumerate}
    \item \textbf{Heave (Vertical Motion):} Achieved by uniformly increasing or decreasing the speeds of all four rotors, which changes the net vertical thrust relative to the force of gravity.
    \item \textbf{Roll Motion ($\phi$):} Achieved by establishing a thrust differential between the lateral rotor pairs. For instance, in an ``X'' configuration, increasing the speeds of the left-hand rotors while decreasing the speeds of the right-hand rotors creates a rolling moment about the longitudinal $x$-axis.
    \item \textbf{Pitch Motion ($\theta$):} Achieved by establishing a thrust differential between the longitudinal rotor pairs. Increasing the speeds of the rear rotors while decreasing the speeds of the front rotors generates a pitching moment about the lateral $y$-axis.
    \item \textbf{Yaw Motion ($\psi$):} Generated by varying the torque balance between the CW and CCW pairs. For example, increasing the speeds of the CCW rotors (1 and 3) while decreasing the speeds of the CW rotors (2 and 4) causes an unbalanced reactive moment, turning the vehicle about its vertical $z$-axis.
\end{enumerate}

Because a quadrotor is an underactuated system—possessing six degrees of freedom but only four independent control inputs—motion along the horizontal translational axes ($x$ and $y$) cannot be commanded directly. Instead, horizontal translation is achieved through a dynamic coupling where the vehicle must tilt its thrust vector by executing roll or pitch rotations, creating a lateral component of force in the desired direction of travel.

\section{The Crazyflie Nano-Quadrotor Ecosystem}
The Crazyflie, developed by Bitcraze AB, is an open-source, ultra-lightweight nano-UAV system designed specifically as a development and testing platform for indoor robotics, control engineering, and swarm flight research \cite{ref_41, ref_40}. By combining a completely open-source firmware codebase with accessible hardware schematics and an array of expansion decks, the platform bridges the gap between theoretical algorithm design and physical experimental validation.

\subsection{Historical Versions and Evolution}
The Crazyflie platform has evolved through multiple iterations, with each version introducing structural, sensing, and processing improvements to meet the requirements of aerial robotics researchers:
\begin{itemize}
    \item \textbf{Crazyflie 1.0 (2013):} The foundational model, which introduced the concept of an open-source development nano-quadrotor. It relied on a lightweight printed circuit board (PCB) as its primary structural frame and brushed coreless DC motors. High-level control and state estimation were highly constrained due to limited processing resources.
    \item \textbf{Crazyflie 2.0 (2014):} Introduced a dual-MCU processing architecture that separated high-level application code from low-level communication and power management. Sensor upgrades included an integrated Inertial Measurement Unit (IMU) with a gyroscope, accelerometer, and magnetometer.
    \item \textbf{Crazyflie 2.1 (2019):} The baseline platform used in this work. It features an upgraded onboard sensor suite, improved radio telemetry performance, a more robust mechanical frame, and an expansion connector that supports automatic detection of expansion decks \cite{ref_40}.
    \item \textbf{Crazyflie 2.1+ (2024):} An incremental update to the Crazyflie 2.1 platform, incorporating an optimized battery chemistry and redesigned propeller geometries, yielding a 15\% improvement in flight performance and thrust generation.
    \item \textbf{Crazyflie 2.1 Brushless (2025):} The latest development in the platform lineup. Retaining the 92 mm motor-to-motor diagonal dimensions of the standard coreless versions, it replaces the brushed DC motors with high-efficiency, high-durability brushless motors, increasing the payload capacity to 40 g and providing a larger thrust-to-weight ratio to support demanding maneuvers \cite{ref_53}.
\end{itemize}

\subsection{Detailed Specifications of the Research Platform}
This research focuses on the Crazyflie 2.1 and its brushless counterparts. The key physical, sensing, and electrical specifications of the platform are summarized in Table \ref{tab:cf_specs_detailed}.

\begin{table}[h!]
\centering
\caption{Crazyflie 2.1 and Brushless Version Physical and Electrical Specifications}
\label{tab:cf_specs_detailed}
\begin{tabular}{|l|c|c|}
\hline
\textbf{Specification} & \textbf{Crazyflie 2.1 (Standard)} & \textbf{Crazyflie 2.1 Brushless} \\
\hline
Takeoff Weight (Standard) & 29 g & 34 g \\
Max Recommended Payload & 15 g & 40 g \\
Diagonal Dimension & 92 mm & 100 mm \\
Primary Processor & STM32F405 (168 MHz) & STM32F405 (168 MHz) \\
Secondary Processor & nRF51822 (32 MHz) & nRF51822 (32 MHz) \\
Inertial Measurement Unit & BMI088 (Accel/Gyro) & BMI088 (Accel/Gyro) \\
Barometer Sensor & BMP388 (Pressure) & BMP388 (Pressure) \\
Battery Capacity & 250 mAh LiPo & 450 mAh LiPo \\
Nominal Flight Time & 7 minutes & 8 to 10 minutes \\
\hline
\end{tabular}
\end{table}

\section{Operational Architecture: How Crazyflie Works}
The operational workflow of the Crazyflie is based on its dual-MCU architecture, which decouples tasks according to their real-time execution constraints \cite{ref_40}. High-frequency flight calculations and hardware interactions are performed on-board, while high-level positioning and mission scripting are conducted off-board.

\subsection{Onboard Dual-MCU Processing}
\begin{enumerate}
    \item \textbf{Low-Level Radio and Power SoC (nRF51822):} Operating at 32 MHz, this ARM Cortex-M0 processor continuously manages the 2.4 GHz low-latency radio link, the Bluetooth Low Energy (BLE) stack, and the onboard battery charging loop. It offloads low-level communication and power monitoring from the primary flight processor, ensuring that high-level flight controls are executed without interruption.
    \item \textbf{Main Flight Control MCU (STM32F405):} A 32-bit ARM Cortex-M4 processor equipped with a floating-point unit (FPU) running at 168 MHz. This processor executes the main flight control loop at 500 Hz. In each execution cycle, the MCU reads the high-frequency accelerometer and gyroscope data from the BMI088 sensor, estimates the vehicle state using an Extended Kalman Filter (EKF), processes the High-Level Commander trajectories, runs the attitude and position control loops, and generates Pulse Width Modulation (PWM) signals to control the motor drivers \cite{ref_54}.
\end{enumerate}

\subsection{State Estimation and Indoor Positioning}
Because a nano-quadrotor's onboard IMU suffers from structural vibrations and sensor drift, achieving autonomous flight requires the integration of high-precision indoor localization systems to supplement the onboard sensors. The Crazyflie fuses these measurements in its onboard EKF \cite{ref_41, ref_54}:
\begin{itemize}
    \item \textbf{Lighthouse Positioning System:} Employs HTC Vive laser base stations that sweep the workspace with infrared beams. The onboard Lighthouse deck decodes the timing of these sweeps to calculate the drone's 3D pose directly on the vehicle, enabling decentralization \cite{ref_55}.
    \item \textbf{Loco Positioning System (UWB):} Uses Decawave DWM1000 Ultra-Wideband (UWB) transceiver nodes to measure distance. The Crazyflie uses a Time Difference of Arrival (TDoA) algorithm to estimate its coordinates in GPS-denied environments \cite{ref_55}.
    \item \textbf{Motion Capture System (MoCap):} Employs external infrared camera systems (e.g., Vicon or OptiTrack) to track retroreflective markers on the drone. The ground-truth state is broadcast to the Crazyflie via a 2.4 GHz Crazyradio PA at 100 Hz, with position errors of less than 1 mm \cite{ref_39}.
\end{itemize}

\section{Rationale for Platform Selection}
Selecting the Crazyflie 2.1 nano-quadrotor as our experimental platform is motivated by key practical advantages:
\begin{enumerate}
    \item \textbf{Exceptional Operational Safety:} Traditional research quadrotors (weighing 1 kg or more) present severe safety hazards to operators and laboratory equipment in the event of controller failure. Weighing only 29 g to 45 g, the Crazyflie has minimal kinetic energy, enabling safe indoor experimentation without requiring protective netting or massive flight volumes \cite{ref_41, ref_46}.
    \item \textbf{High Control Agility:} The low rotational inertia of the nano-scale frame allows the Crazyflie to execute aggressive, high-bandwidth maneuvers, making it an excellent platform for testing nonlinear robust control algorithms.
    \item \textbf{Extensible Expansion Bus:} The expansion interface exposes I2C, SPI, UART, and GPIO lines, enabling the attachment of additional decks (such as the Flow Deck v2 or the Multi-Ranger) to support obstacle avoidance and optical flow state estimation without requiring physical rewiring or soldering \cite{ref_40}.
    \item \textbf{Robust Software Ecosystem:} The open-source firmware facilitates high-frequency logging and parameter adjustments over the air. Integration with modern frameworks like Crazyswarm2 and ROS 2 Humble allows for seamless simulation-to-reality transition \cite{ref_39}.
\end{enumerate}

\section{Control System Requirements and Challenges}
A quadrotor is an open-loop unstable dynamical system. Without continuous feedback control, even minute asymmetries in motor thrust, aerodynamic disturbances, or sensor noise will cause the vehicle to diverge from steady flight and crash. Developing robust control loops for the Crazyflie platform must address several operational challenges:

\subsection{Aerodynamic Drag and Unmodeled Dynamics}
At high translational velocities or when operating in turbulent environments, the quadrotor is subject to nonlinear rotor drag, blade flapping, and ground effect forces. These disturbances are highly nonlinear and difficult to model analytically from first principles, appearing as unstructured model uncertainties in the dynamic equations.

\subsection{Strict Payload Constraints}
The maximum recommended payload for the coreless Crazyflie 2.1 is strictly limited to 15 g. This prevents the integration of high-performance companion computers (such as the NVIDIA Jetson series) directly onto the vehicle. High-order, computationally intensive control laws must therefore be designed to run efficiently on the 168 MHz STM32F405 microcontroller.

\subsection{Battery Voltage Decay (Battery Sag)}
The flight of the Crazyflie relies on a small single-cell (1S) Lithium-Polymer battery. Under full propeller load, the battery open-circuit voltage decays from 4.1 V to its nominal 3.7 V during flight, as shown conceptually in Figure \ref{fig:battery_sag_response}. This sag directly reduces the motor thrust coefficient $K_f$, meaning the same PWM command generates less thrust over time \cite{ref_40}. A robust controller must actively compensate for this time-varying parametric drift to maintain altitude tracking accuracy.

\begin{figure}[h!]
\centering
\fbox{
  \begin{minipage}{0.8\textwidth}
    \centering
    \vspace{1.5cm}
    \textbf{}
    \vspace{1.5cm}
  \end{minipage}
}
\caption{Conceptual representation of battery voltage sag under propeller load during standard hover flight, illustrating the decay of effective thrust.}
\label{fig:battery_sag_response}
\end{figure}

\section{Mathematical Formulation Framework}
To establish the control laws, we must define the mathematical variables that characterize the quadrotor's state and actuator inputs. The quadrotor is modeled as a rigid body with mass $m$ and a diagonal inertia matrix $J \in \mathbb{R}^{3 \times 3}$.

We define the state vector $X(t) \in \mathbb{R}^{12}$ in terms of position, velocity, orientation, and body-fixed angular rates:
\begin{equation}
X(t) = \begin{bmatrix} p(t) \\ v(t) \\ \eta(t) \\ \omega(t) \end{bmatrix} \in \mathbb{R}^{12}
\end{equation}
where:
\begin{itemize}
    \item $p(t) = [x(t), y(t), z(t)]^T \in \mathbb{R}^3$ is the linear position of the quadrotor's center of mass relative to the Inertial World Frame ($W$).
    \item $v(t) = [\dot{x}(t), \dot{y}(t), \dot{z}(t)]^T \in \mathbb{R}^3$ is the linear velocity vector in the World Frame.
    \item $\eta(t) = [\phi(t), \theta(t), \psi(t)]^T \in \mathbb{R}^3$ represents the Euler orientation angles (roll, pitch, yaw) mapping the Body Frame ($B$) to the World Frame ($W$).
    \item $w(t) = [p_w(t), q_w(t), r_w(t)]^T \in \mathbb{R}^3$ contains the angular velocities of the vehicle resolved in the Body-fixed Frame.
\end{itemize}

The control input vector $U(t) \in \mathbb{R}^4$ comprises the virtual forces and torques generated by the four propellers:
\begin{equation}
U(t) = \begin{bmatrix} U_1(t) \\ U_2(t) \\ U_3(t) \\ U_4(t) \end{bmatrix} = \begin{bmatrix} T(t) \\ \tau_{\phi}(t) \\ \tau_{\theta}(t) \\ \tau_{\psi}(t) \end{bmatrix}
\end{equation}
where $U_1 = T$ is the net vertical thrust along the body $z$-axis, and $U_2 = \tau_{\phi}$, $U_3 = \tau_{\theta}$, and $U_4 = \tau_{\psi}$ represent the stabilizing torques about the body-fixed $x$, $y$, and $z$ axes, respectively. 

This mathematical framework underpins the control problem addressed in this thesis: synthesizing a robust, non-singular nonlinear control law that maps the tracking errors of the state vector $X(t)$ to the control input vector $U(t)$ under parametric uncertainties and external wind disturbances. The exact kinematic and dynamic equations mapping the actuator inputs to these twelve states will be derived from first principles using Newton-Euler formulations in Chapter \ref{ch:dynamics}.

% physics of quadrotors \\
% about crazyflie \\
% its versions \\
% current version of crazyflie \\
% how crazyflie works \\
% reason of choosing crazyflie \\
% crazyflie images \\
% crazyflie citations \\
% need of controls here \\
% why robust controls now? \\
% how to describe the system mathematically to control it? \\